\documentclass[../../../uem_phy_std_q.tex]{subfiles}

\begin{document}
    図のように，真空中において，同形の3枚の金属極板X，Y，Zを%
    平行に狭い間隔で向き合わせてあり，電池$\mathrm{B}_1$（起電力$2V_0$），$\mathrm{B}_2$（起電力$V_0$），%
    スイッチ$\mathrm{S}_1$，$\mathrm{S}_2$およびアースと接続されて回路を形成している．%
    XとYの間隔は$d$，YとZの間隔は$2d$である．%
    極板の面積を$S$，真空の誘電率を$\varepsilon_0$とし，端の効果は無視する．%
    はじめ，$\mathrm{S}_1$，$\mathrm{S}_2$は共に開いており，いずれの極板も帯電していない．%

    \begin{enumerate}
        \item $\mathrm{S}_1$を閉じて安定するまで待った．%
            このとき，X下面の電荷$x_1$，Y下面の電荷$y_1$を求めよ．
        \item $\mathrm{S}_1$を開き，$\mathrm{S}_2$を閉じて安定するまで待った．%
            このとき，X下面の電荷$x_2$，Y下面の電荷$y_2$を求めよ．
        \item $\mathrm{S}_2$を開き，$\mathrm{S}_1$を閉じて安定するまで待った．%
            このとき，X下面の電荷$x_3$，Y下面の電荷$y_3$を求めよ．
        \item スイッチはそのままで，XとYの間を誘電体（比誘電率$\varepsilon_r$）で満たし，%
            安定するまで待った．このとき，X下面の電荷$x_4$，Y下面の電荷$y_4$を求めよ．
    \end{enumerate}
    \begin{minipage}{\linewidth}
        \centering
            \includegraphics{\subfix{fig_uem_phy_std_em_cb_11/fig_uem_phy_std_em_cb_11_00_q.pdf}}
    \end{minipage}
\end{document}




